%%%%%%%%%%%%%%%%%%%%%%%%%%%%%%%%%%%%%%%%%%%%%%%%%%%%%%%%%%%%%
%                                                           %
% DATE: 30 - July - 2001                                    %
%                                                           %
% Written with AMS-LaTeX 2.0                                %
%                                                           %
% For problems with this file, write to                     %
%    Juan L. Varona                                         %
%    Dept. Mathematics and Computation                      %
%    University of La Rioja, SPAIN                          %
%    jvarona@dmc.unirioja.es                                %
%    http://www.unirioja.es/dptos/dmc/jvarona/welcome.html  %
%%%%%%%%%%%%%%%%%%%%%%%%%%%%%%%%%%%%%%%%%%%%%%%%%%%%%%%%%%%%%

%%%%%%%%%%%%%%%%%%%%%%%%%%%%%%%%%%%%%%%%%%%%%%%%%%%%%%%%%%%%%
% Beginning of the paper                                    %
%%%%%%%%%%%%%%%%%%%%%%%%%%%%%%%%%%%%%%%%%%%%%%%%%%%%%%%%%%%%%
%
% AMS-LaTeX 2.0 file for journals:
%
%%%%%\documentclass{amsart} %%%%% Generic

\documentclass{amsart}

%% \renewcommand{\baselinestretch}{1.79} % NEAR DOUBLE SPACE
% DOBLE ESPACIO ES CON "2".
% CON "1.8" YA ME SALEN 20 PAGINAS.
% CON "2", ME SALEN 22; Y, ADEMAS, LA PRIMERA QUEDA HORRIBLE.

%\date{\today}
%\date{July 30, 2001}

%\usepackage[notcite,notref]{showkeys}

% Lemmas, theorems and corolaries:
\newtheorem{lem}{Lemma}
\newtheorem{thm}{Theorem}

\theoremstyle{remark}
\newtheorem{rem}{Remark}

\newcommand{\N}{\mathbb{N}}
\newcommand{\R}{\mathbb{R}}
\renewcommand{\Re}{\operatorname{Re}}
\newcommand{\iden}{\operatorname{Id}}

% Optional argument -> \delta
% Usually, we use \Jad without argument.
% To put something optional "opt", use \Jad[opt]
\newcommand{\Jad}[1][\delta]{\mathcal{J}_{\alpha}^{#1}}

%\newcommand{\htr}[1]{\mathcal{H}_{#1}}
%\newcommand{\htrk}[2]{\mathcal{H}_{#1}^{#2}}
\newcommand{\htr}{\mathcal{H}_{\alpha}}
\newcommand{\htrk}{\mathcal{H}_{\alpha}^{k}}

\newcommand{\ma}{M_{\alpha}}
\newcommand{\mak}{M_{\alpha,k}}

\newcommand{\BR}{M_\alpha^{\delta}}
\newcommand{\BRk}{M_{\alpha,k}^{\delta}}
\newcommand{\BRkd}[1]{M_{\alpha,k}^{\delta(#1)}} % Analytic family

\newcommand{\KBRk}{K_{\alpha+k}^{\delta}} % The kernel
\newcommand{\KBRkd}[1]{K_{\alpha+k}^{\delta(#1)}} % Analytic family

% Optional argument -> s
% Usually, we use \calT without argument.
% To put something optional \calT[opt]
\newcommand{\calT}[1][s]{\mathcal{T}^{#1}_{\alpha,k}}

% Optional argument -> \alpha+k
\newcommand{\calK}[1][\alpha+k]{\mathcal{K}^s_{#1}}

% L^p(x^{2\alpha+1})-norm of a function:
\newcommand{\lpnorm}[2]{\|{#1}\|_{L^{#2}(x^{2\alpha+1})}}
%\newcommand{\lpnormt}[2]{\|{#1}\|_{L^{#2}(t^{2\alpha+1})}}

\newcommand{\biglpnorm}[2]{\left\|{#1}\right\|_{L^{#2}(x^{2\alpha+1})}}

%%%%%%%%%%%%%%%%%%%%%%%%%%%%%%%%%%%%%%%%%%%%%%%%%%%%%%%%%%%%%%%%%%%%%%%

\begin{document}

\title[Bochner-Riesz operators]{An uniform boundedness
for Bochner-Riesz operators related to the Hankel transform}

\thanks{Research supported by grants
BFM2000-0206-C04-03 of the DGI and API-01/B38 of the UR}

\author[\'O. Ciaurri]{\'Oscar Ciaurri}
\address{Departamento de Matem\'aticas y Computaci\'on,
Universidad de La Rioja, Edificio J.~L.~Vives,
Calle Luis de Ulloa s/n, 26004 Logro\~no, Spain}
\email{oscar.ciaurri@dmc.unirioja.es}

\author[J. L. Varona]{Juan L. Varona}
\address{Departamento de Matem\'aticas y Computaci\'on,
Universidad de La Rioja, Edificio J.~L.~Vives,
Calle Luis de Ulloa s/n, 26004 Logro\~no, Spain}
\email{jvarona@dmc.unirioja.es}
\urladdr{http://www.unirioja.es/dptos/dmc/jvarona/welcome.html}

% This need amsart version 2.0 or later
% (remove "[2000]" for amsart version 1.2)
\subjclass[2000]{Primary 42A45; Secondary 42A38}

\keywords{Bochner-Riesz operator, multipliers, Hankel transform,
Fourier-Neumann series.}

\thanks{PUBLISHED IN: J.\ Inequal.\ Appl.\ \textbf{7} (2002), no.~6, 759--777.}

% SUELE SER UNA BUENA IDEA NO USAR DEFINICIONES DE TeX EN EL ABSTRACT,
% PARA QUE LO PUEDAN COPIAR Y PEGAR A DISTINTOS SITIOS.

%%%%%%%%%%%%%%%%%%%%%%%%%%%%%%%%%%%%%%%%%%%%%%%%%%%%%%%%%%%%%%%%%%%%%%%%%%%
\begin{abstract}
Let $\mathcal{H}_{\alpha}$ be the modified Hankel transform
\begin{equation*}
    \mathcal{H}_{\alpha}(f,x) = \int_0^\infty
       \frac{J_\alpha(xt)}{(xt)^\alpha} \, f(t) t^{2\alpha+1} \, dt,
\end{equation*}
defined for suitable functions and extended to some
$L^p((0,\infty),x^{2\alpha+1})$ spaces. Given $\delta>0$, let
$M_{\alpha}^{\delta}$ be the Bochner-Riesz operator for the Hankel
transform. Also, we take the following generalization
\begin{equation*}
    \mathcal{H}_{\alpha}^k(f,x) = \int_0^\infty
        \frac{J_{\alpha+k}(xt)}{(xt)^\alpha} \, f(t) t^{2\alpha+1} \, dt,
    \quad k = 0,1,2,\dots
\end{equation*}
for the Hankel transform, and define $M_{\alpha,k}^{\delta}$ as
\begin{equation*}
    M_{\alpha,k}^{\delta} f = \mathcal{H}_{\alpha}^k \left(
       (1-x^2)^{\delta}_+ \mathcal{H}_{\alpha}^k f \right),
    \quad k = 0,1,2,\dots
\end{equation*}
(thus, in particular, $M_{\alpha}^{\delta} = M_{\alpha,0}^{\delta}$).
In the paper, we study the uniform boundedness of
$\{M_{\alpha,k}^{\delta}\}_{k\in\N}$
in $L^p((0,\infty),x^{2\alpha+1})$ spaces when $\alpha \ge 0$.
We found that, for $\delta > (2\alpha+1)/2$ (the critical index), the uniform
boundedness of $\{M_{\alpha,k}^{\delta}\}_{k=0}^\infty$ is satisfied for every
$p$ in the range $1 \le p \le \infty$. And, for $0 < \delta \le (2\alpha+1)/2$, the
uniform boundedness happens if and only if
\begin{equation*}
    \frac{4(\alpha+1)}{2\alpha+3+2\delta} < p
    < \frac{4(\alpha+1)}{2\alpha+1-2\delta}.
\end{equation*}
In the paper, the case $\delta = 0$ (the corresponding generalization of the
$\chi_{[0,1]}$-multiplier for the Hankel transform) is previously analyzed;
here, for $\alpha > -1$. For this value of $\delta$, the uniform boundedness
of $\{M_{\alpha,k}^0\}_{k=0}^\infty$ is related to the convergence of
Fourier-Neumann series.
\end{abstract}
%%%%%%%%%%%%%%%%%%%%%%%%%%%%%%%%%%%%%%%%%%%%%%%%%%%%%%%%%%%%%%%%%%%%%%%%%%%

%%%%%%%%%%%%%%%%%%%%%%%%%%%%%%%%%%%%%%%%%%%%%%%%%%%%%%%%%%%%%%%%%%%%%%%%%%%
\maketitle
%%%%%%%%%%%%%%%%%%%%%%%%%%%%%%%%%%%%%%%%%%%%%%%%%%%%%%%%%%%%%%%%%%%%%%%%%%%

%%%%%%%%%%%%%%%%%%%%%%%%%%%%%%%%%%%%%%%%%%%%%%%%%%%%%%%%%%%%%%%%%%%%%%%%%%%
\section{Introduction}
%%%%%%%%%%%%%%%%%%%%%%%%%%%%%%%%%%%%%%%%%%%%%%%%%%%%%%%%%%%%%%%%%%%%%%%%%%%

Let $\alpha\ge-1/2$. For a function $f(t)$ on the interval $(0,\infty)$,
the so-called modified Hankel transform $\htr(f,x)$, $x>0$, of order $\alpha$
is given by
\begin{equation}
\label{ec:htr}
    \htr(f,x) = \int_0^\infty \frac{J_\alpha(xt)}{(xt)^\alpha} \, f(t)
      t^{2\alpha+1} \, dt,
\end{equation}
where $J_\alpha(x)$ is the Bessel function of the first kind of order~$\alpha$.
Well-known bounds for Bessel function and H\"older's inequality
show that~(\ref{ec:htr}) is well defined for every
$f\in L^p((0,\infty),x^{2\alpha+1}\,dx)$
($L^p(x^{2\alpha+1})$, from now on) with
$1 \le p < \frac{4(\alpha+1)}{(2\alpha+3)}$.

Furthermore, it is easy to see that $\htr$ is a bounded
operator from $L^1(x^{2\alpha+1})$ into $L^\infty(x^{2\alpha+1})$. Also, as
usual, the expression~(\ref{ec:htr}) is extended by continuity to different
$L^p(x^{2\alpha+1})$ spaces. For instance, it is well known that $\htr$ is an
isomorphism from $L^2(x^{2\alpha+1})$ into itself and $\htr \circ \htr =
\iden$ (for $\alpha>-1$, the extension of $\htr$ to $L^2(x^{2\alpha+1})$ can
also be done; see~\cite{Dur1,Dur2,Var,BeStem}). As a consequence, of course,
we can get the corresponding interpolation result that we do not detail here.

Associated with $\htr$, we can define the ball multiplier $\ma$ as the
operator that verifies the relation $\htr(\ma f) = \chi_{[0,1]}\htr f$; or,
in other words,
\begin{equation*}
    \ma f = \htr(\chi_{[0,1]}\htr f).
\end{equation*}
This operator is bounded from $L^p(x^{2\alpha+1})$ into itself if and only
if $\frac{4(\alpha+1)}{2\alpha+3} < p < \frac{4(\alpha+1)}{2\alpha+1}$.
The first proof of this fact can be found in~\cite{Her}. For this and related
properties on Hankel multipliers see
also~\cite{Iga,GasTr,Var,StemTr,GPV-Con,CGPV-FN,CGPV-DIE}.

Now, let us take the following generalization for the Hankel transform. We
consider
\begin{equation}
\label{ec:htrk}
    \htrk(f,x) = \int_0^\infty \frac{J_{\alpha+k}(xt)}{(xt)^\alpha} \, f(t)
      t^{2\alpha+1} \, dt,
    \quad k = 0,1,2,\dots.
\end{equation}
In this way, the operator $\ma$ can be also generalized by taking
\begin{equation}
\label{ec:mak}
    \mak f = \htrk(\chi_{[0,1]}\htrk f),
    \quad k = 0,1,2,\dots.
\end{equation}
The study of the uniform boundedness of these operators in
$L^p(x^{2\alpha+1})$ spaces is very useful. As we can see
in~\cite{Var,GPV-Con,CGPV-FN}, to prove the mean convergence
of Fourier-Neumann series is reduced to prove the uniform boundedness of
the operators $\{\mak f\}_{k\in\N}$. There, it is proved
that the uniform boundedness in $L^p(x^{2\alpha+1})$ is equivalent to
$\max\{\frac43,\frac{4(\alpha+1)}{2\alpha+3}\} < p <
\min\{4,\frac{4(\alpha+1)}{2\alpha+1}\}$.
Actually, in this article, we will explicitly show the uniform
boundedness of $\{\mak f\}_{k\in\N}$ only because it plays an important role
in the proof of some other results of the paper. Moreover, the boundedness for
$\mak$ is shown not only for $\alpha\ge-1/2$, but also for $\alpha>-1$.

Another multiplier for the Hankel transform is the Bochner-Riesz multiplier;
of course, it is similar to the well known Bochner-Riesz multiplier for the
Fourier transform (see, for instance,~\cite{Well,RF} or~\cite[Ch.~IX,
\S~2.2]{St-HA}). Taking $\delta>0$, the operator $\BR$ is the one that
makes $\htr(\BR f) = (1-x^2)^{\delta}_+ \htr f$,
being $(1-x^2)_+ = \max\{0,1-x^2\}$ (of course, $\BR$ for $\delta=0$ would be
$\ma$). Also, it can be described as
\begin{equation}
\label{ec:BR}
    \BR f = \htr \left( (1-x^2)^{\delta}_+ \htr f \right).
\end{equation}
Again, Bochner-Riesz multiplier can be generalized
by using~(\ref{ec:htrk}). Thus, we take
\begin{equation}
\label{ec:BRk}
    \BRk f = \htrk \left( (1-x^2)^{\delta}_+ \htrk f \right),
    \quad k = 0,1,2,\dots.
\end{equation}
Similarly to the ball multiplier generalization, we can wonder if these
operators are uniformly bounded in $L^p(x^{2\alpha+1})$. This paper is
devoted to the study of this fact. We will only deal with the case $\alpha \ge
0$.

We found that there exists an index $\delta_0 = (2\alpha+1)/2$
such that, for $\delta > \delta_0$, the uniform boundedness is true
for every $p$ in the range $1 \le p \le \infty$; 
however, for $0 < \delta \le \delta_0$, the uniform boundedness
only happens in a finite range of~$p$.
Moreover, as we will see, the proofs for both facts are different.
This value $\delta_0$ is called the critical index. The existence of a
critical index is also a common fact in the study of the boundedness of the
Bochner-Riesz operator for the Fourier transform (see, for
instance,~\cite{StWe})



The paper is organized as follows. 
In Section~\ref{sec:main}, we give the main results of the paper. First, the
uniform boundedness for $\{\mak\}_{k\in\N}$ (Theorem~\ref{thm:mak}); then,
the uniform boundedness for $\{\BRk\}_{k\in\N}$, both for $\delta > \delta_0$
(Theorem~\ref{thm:encima}) and for $0 < \delta \le \delta_0$
(Theorem~\ref{thm:debajo}).
In Section~\ref{sec:mak}, we give the proof of Theorem~\ref{thm:mak}, which is
reduced to known results that already appear in papers related with the
uniform boundedness of Fourier-Neumann series.
In Section~\ref{sec:encima} we give the proof of Theorem~\ref{thm:encima}.
To obtain it, we will use some results about the
translation for the Hankel transform and the corresponding convolution
product. Finally, in Section~\ref{sec:debajo}, we give the proof of
Theorem~\ref{thm:debajo}. In this case, we will apply an interpolation result
for families of analytic operators from Stein~\cite{St-Int}, and the uniform
boundedness of $\{\mak\}_{k\in\N}$ that appears in Theorem~\ref{thm:mak}.

Throughout this paper, we will use $C$ (or $C'$) to denote a positive constant
independent of $f$ and $k$ (and of any other variable, if it is the case),
which can assume different values in different occurrences. When $C$ has a
subindex, it depends only on the parameters that appear in the subindex.

Some of the operators defined in this paper have an integral expression only
for suitable functions, and then they are extended by density and continuity.
Usually, a class $S^+$ of smooth functions in $(0,\infty)$, related with
Schwartz class~$S$, is used ($S^+$ are the even functions of $S$ restricted to
$(0,\infty)$); see, for instance,~\cite{Dur2,ColCTV} for details
(and~\cite{GasTr,StemTr} for related density results). Moreover, with these
functions, Fubini's Theorem can be applied when necessary. This is a standard
technique; we will implicitly use this kind of arguments sometimes without
notice it again.

\begin{rem}
\label{rem:raiz}
Note that, although we have referred to some results that appeared
in~\cite{Var,GPV-Con,CGPV-FN, CGPV-DIE}, the Hankel transform used in
these papers (as well as by other authors such as~\cite{Dur1,Dur2}) was
somewhat different. In these papers, instead of~(\ref{ec:htr}), it was used
\begin{equation}
\label{ec:htrRaiz}
    \htr(f,x) = \frac{x^{-\alpha/2}}2
      \int_0^\infty J_{\alpha}(\sqrt{xt}) \, g(t) t^\alpha \, dt,
\end{equation}
defined in $L^p(x^\alpha)$. And the ball multiplier $\ma$ (and $\mak$) was also
defined for this Hankel transform. It is clear that, by the changes of
variables $t \mapsto t^2$ and $x \mapsto x^2$, the Hankel
transform~(\ref{ec:htrRaiz}) in $L^p(x^\alpha)$
becomes~(\ref{ec:htr}) in $L^p(x^{2\alpha+1})$. Of course, the
range of $p$ for which there exists boundedness is preserved.

Nowadays, it seems that the notation in~(\ref{ec:htr}) is more used (see,
for instance~\cite{Hir,Iga,GasTr,StemTr,Men,ColCTV,BeRo,BeStem}, although it is
sometimes called Fourier-Bessel transform), so we have adopted it in this
paper.

When studying the boundedness of $\ma$ and $\mak$, the notation
in~(\ref{ec:htrRaiz}) was more handy than~(\ref{ec:htr}): the
operators that appear are, directly, Hilbert transforms with weights
(with~(\ref{ec:htr}), in the denominator, an $x^2-t^2$ arises instead of a
$x-t$, and this requires extra work). However, in this paper, is more
suitable to use~(\ref{ec:htr}) and its corresponding $\BR$ and~$\BRk$.
\end{rem}


%%%%%%%%%%%%%%%%%%%%%%%%%%%%%%%%%%%%%%%%%%%%%%%%%%%%%%%%%%%%%%%%%%%%%%%%%%%
\section{Main results}
\label{sec:main}
%%%%%%%%%%%%%%%%%%%%%%%%%%%%%%%%%%%%%%%%%%%%%%%%%%%%%%%%%%%%%%%%%%%%%%%%%%%

Let us begin to show an expression for the $\chi_{[0,1]}$-multiplier
generalization $\mak f$ described in~(\ref{ec:mak}); or, in other words,
$\BRk f$ for $\delta=0$. We will study $\{\mak\}_{k\in\N}$ not only for
$\alpha\ge-1/2$, but also for $\alpha>-1$. By using~(\ref{ec:mak}) and
Fubini's Theorem we obtain
\begin{equation}
\label{ec:makConKernel}
   \mak(f,x) = x^{-\alpha} \int_0^\infty
     \left( \int_0^1 J_{\alpha+k} (ts) J_{\alpha+k} (xs) s \, ds \right)
       t^{\alpha+1} f(t) \, dt.
\end{equation}

The result that shows the uniform boundedness of this family of operators
is the following:

\begin{thm}
\label{thm:mak}
Let $\alpha > -1$, $1 < p < \infty$, and the family of operators
$\{\mak\}_{k\in\N}$ defined as in~(\ref{ec:mak}). Then,
\begin{equation*}
%\label{ec:makAcot}
    \lpnorm{\mak f}{p} \le C \lpnorm{f}{p},
    \quad k\in \N,
\end{equation*}
if and only if
\begin{equation*}
\begin{cases}
    4/3 < p < 4, &\text{if \ $-1 < \alpha < 0$,} \\
    \frac{4(\alpha+1)}{2\alpha+3} < p 
      < \frac{4(\alpha+1)}{2\alpha+1}, &\text{if \ $0 \le \alpha$.}
\end{cases}
\end{equation*}
\end{thm}

Now, let us consider the generalized Bochner-Riesz multiplier $\BRk$ described
in~(\ref{ec:BRk}). Let us take $\delta>0$.
An expression that allows us to show $\BRk$ as an operator with a kernel
$\KBRk$ is the following:
\begin{equation*}
%\label{ec:BRkConKernel}
    \BRk(f,x)
    = x^{-\alpha} \int_0^{\infty} f(t) \KBRk(x,t) t^{\alpha+1} \, dt
\end{equation*}
with
\begin{equation}
\label{ec:KernelBRkintermedia}
    \KBRk(x,t)
    = \int_0^1 s (1-s^2)^{\delta}
      J_{\alpha+k}(xs)J_{\alpha+k}(ts) \, ds.
\end{equation}
Again, it can be easily deduced by applying Fubini's Theorem to the
definition for~$\BRk$ given in~(\ref{ec:BRk}).

The main results in the paper are the ones that show uniform boundedness for
$\{\BRk\}_{k\in\N}$. When $\delta$ exceeds the critical index, we have

\begin{thm}
\label{thm:encima}
Let $\alpha\ge 0$, $\delta > \delta_0 = (2\alpha+1)/2 $ and $1 \le p \le
\infty$, and the family of operators $\{\BRk\}_{k\in\N}$ defined as
in~(\ref{ec:BRk}). Then,
\begin{equation*}
%\label{ec:BRkencima}
    \lpnorm{\BRk f}{p}
    \le C \lpnorm{f}{p},
    \quad \ k\in \N.
\end{equation*}
\end{thm}

And, for $\delta$ below the critical index, we have

\begin{thm}
\label{thm:debajo}
Let $\alpha\ge 0$, $0 < \delta \le \delta_0 = (2\alpha+1)/2$ and $1 \le p \le
\infty$, and the family of operators $\{\BRk\}_{k\in\N}$ defined as
in~(\ref{ec:BRk}). Then,
\begin{equation*}
%\label{ec:BRkdebajo}
    \lpnorm{\BRk f}{p} \le C \lpnorm{f}{p},
    \quad k\in \N,
\end{equation*}
if and only if
\begin{equation}
\label{ec:condprin}
    \frac{4(\alpha+1)}{2\alpha+3+2\delta}
    < p <
    \frac{4(\alpha+1)}{2\alpha+1-2\delta}.
\end{equation}
\end{thm}

%%%%%%%%%%%%%%%%%%%%%%%%%%%%%%%%%%%%%%%%%%%%%%%%%%%%%%%%%%%%%%%%%%%%%%%%%%%
\section{Proof of Theorem~\ref{thm:mak} (i.e., uniform boundedness
for $\delta=0$)}
\label{sec:mak}
%%%%%%%%%%%%%%%%%%%%%%%%%%%%%%%%%%%%%%%%%%%%%%%%%%%%%%%%%%%%%%%%%%%%%%%%%%%

First, we are going to establish some new expressions
for~(\ref{ec:makConKernel}) that will be more useful to study the uniform
boundedness of the operators.

We will use von Lommel's formula
\begin{align*}
   \int_0^1 J_\nu(ts) J_\nu(xs) s \, ds
   &= \frac{1}{x^2-t^2}
     \left( J_\nu (x) t J'_\nu (t) - J_\nu(t) x J'_\nu(x) \right) \\
   &= \frac{1}{x^2-t^2}
     \left( x J_\nu(t) J_{\nu+1}(x) - t J_{\nu+1}(t) J_\nu (x) \right)
\end{align*}
(for the last equality, use $zJ'_\nu(z) = \nu J_\nu(z) - zJ_{\nu+1}(z)$).
By applying it to~(\ref{ec:makConKernel}) we get
\begin{align}
\label{ec:makJJ}
    \mak(f,x) &=
    x^{-\alpha+1} J_{\alpha+k+1}(x)
        \int_0^\infty \frac{t^{-\alpha} J_{\alpha+k}(t)}{x^2-t^2}
          f(t) t^{2\alpha+1} \, dt \\
    &\qquad - x^{-\alpha} J_{\alpha+k}(x)
        \int_0^\infty \frac{t^{-\alpha+1} J_{\alpha+k+1}(t)}{x^2-t^2}
          f(t) t^{2\alpha+1} \, dt \notag \\
    &= W_{1,k}(f,x) - W_{2,k}(f,x) \notag
\end{align}
or, also,
\begin{align}
\label{ec:makJJ'}
    \mak(f,x) &=
    x^{-\alpha} J_{\alpha+k}(x)
        \int_0^\infty \frac{t^{-\alpha+1} J'_{\alpha+k}(t)}{x^2-t^2}
          f(t) t^{2\alpha+1} \, dt \\
    &\qquad - x^{-\alpha+1} J'_{\alpha+k}(x)
        \int_0^\infty \frac{t^{-\alpha} J_{\alpha+k}(t)}{x^2-t^2}
          f(t) t^{2\alpha+1} \, dt \notag \\
    &= \widetilde W_{1,k}(f,x) - \widetilde W_{2,k}(f,x). \notag
\end{align}


Now, we have all the necessary for

\begin{proof}[Proof of Theorem~\ref{thm:mak}]

These are the operators that appear in the decomposition of the partial sums
of Fourier-Neumann series to prove their uniform boundedness in
$L^p(x^{2\alpha+1})$ spaces, such as it is studied
in~\cite{Var,GPV-Con,CGPV-FN}. Actually, with this notation, the partial sums
are
$S_nf = W_{1,0}f - W_{2,0}f + \widetilde W_{2,2n+2}f - \widetilde
W_{1,2n+2}f$ (take into account that the notation in these papers is a little
different and, also, that a change of variable for $x$ and $t$ is being used,
as described in Remark~\ref{rem:raiz}).

In this way, a sketch of the proof for the uniform boundedness of $\mak$ is as
follows:

For $k=0$, the decomposition~(\ref{ec:makJJ}) is used. Then, well-known bounds
for $|J_{\alpha}|$ and $|J_{\alpha+1}|$ are applied, and so the proof of the
boundedness of $W_{2,0}$ and $W_{1,0}$ is reduced to the boundedness of the
Hilbert transform with weights. We get that this boundedness is true for
$1 < p < \infty$ when $-1 < \alpha < -1/2$; and for
$\frac{4(\alpha+1)}{2\alpha+3} < p < \frac{4(\alpha+1)}{2\alpha+1}$ when
$\alpha \ge -1/2$. (The case $\alpha<-1/2$ is studied only
in~\cite{GPV-Con}.)

For $k>0$, the decomposition~(\ref{ec:makJJ'}) is used. Now, suitable bounds
for $|J_{\alpha+k}|$ and $|J'_{\alpha+k}|$ are applied. Then,
$\widetilde W_{1,k}$ and $\widetilde W_{2,k}$ are bounded by Hilbert transforms
with weights (that depend of $k$). Finally, uniform $A_p$ weights theory is
used to find the uniform boundedness of these Hilbert transforms.
Here, the condition $4/3 < p < 4$ appears.
\end{proof}


\begin{rem}
As we have commented in the previous proof, the condition $4/3 < p < 4$ does 
not arise if we only analyze the case $k=0$. Then, for the boundedness of
the $\chi_{[0,1]}$-multiplier $\ma$ in $L^p(x^{2\alpha+1})$, $1<p<\infty$ and
$\alpha>-1$, we have
\begin{equation*}
    \lpnorm{\ma f}{p} \le C \lpnorm{f}{p}
    \iff
    \begin{cases}
    1 < p < \infty, &\text{if \ $-1 < \alpha < \frac{-1}2$,} \\
    \frac{4(\alpha+1)}{2\alpha+3} < p 
      < \frac{4(\alpha+1)}{2\alpha+1}, &\text{if \ $\frac{-1}2 \le \alpha$.}
    \end{cases}
\end{equation*}
This result is already implicit in~\cite{GPV-Con}, but not explicitely
stated. And, changing $L^p(x^{2\alpha+1})$ by some other weighted $L^p$ spaces,
also~\cite{Var-Evora} can be seen.

A completely different proof of this fact, based on multipliers for
Fourier-Bessel expansions and transplantion theorems, can be found
in~\cite{BeStem}.
\end{rem}

%%%%%%%%%%%%%%%%%%%%%%%%%%%%%%%%%%%%%%%%%%%%%%%%%%%%%%%%%%%%%%%%%%%%%%%%%%%
\section{Uniform boundedness for $\delta$ over the critical index}
\label{sec:encima}
%%%%%%%%%%%%%%%%%%%%%%%%%%%%%%%%%%%%%%%%%%%%%%%%%%%%%%%%%%%%%%%%%%%%%%%%%%%

First, let us describe the translation and convolution for the Hankel transform
as well as some of its properties. A wider exposition of these results can be
found in~\cite{Hir,Hai,Cho}. Mostly in this section, we can assume
$\alpha>-1/2$, although the proof of Theorem~\ref{thm:encima} will require
$\alpha \ge 0$ (because~(\ref{ec:cotaUltraesf}) is not true for~$\alpha<0$).

We consider the translation operator $T^x$, with
$x \ge 0$, defined, for suitable functions $h$, by
\begin{equation*}
%\label{ec:tras}
    T^x(h,t) = \frac{\Gamma(\alpha+1)}{\sqrt{\pi}\,\Gamma(\alpha+1/2)}
      \int_0^{\pi} h(w)\sin^{2\alpha}{\theta} \, d\theta,
\end{equation*}
where $w^2 = x^2+t^2-2xt\cos\theta$, $x,t \ge 0$.
Using this translation, the convolution operator can be described as
\begin{equation}
\label{ec:convo}
    h*g(x) = \int_0^{\infty} g(t) T^x(h,t) t^{2\alpha+1} \, dt.
\end{equation}
It is not difficult to check that, for $\alpha>-1/2$,
\begin{equation}
\label{ec:acotconvo}
    \lpnorm{h*g}{p} \le \lpnorm{h}{1} \lpnorm{g}{p}
\end{equation}
for $1 \le p \le \infty$.
The convolution structure and its boundedness will play a significant role
in the proof of Theorem~\ref{thm:encima}.

Let us introduce some other notation.
Let $\{P^{(\alpha)}_k\}_{k=0}^\infty$ be the ultraspherical polynomials
of order $\alpha$ (also known as Gegenbauer polynomials). Here, and in the
discussion that will follow, the usual modifications must be applied if
$\alpha=0$ (i.e., the use of Chebyshev polynomials). A wide information on
ultraspherical polynomials can be found in~\cite{Sze}. To simplify some
expressions, we take
\begin{equation}
\label{ec:ultraesf}
    p^{\alpha}_k(x)
    = \frac{2^{\alpha-1}k!\,\Gamma(\alpha)}{\pi \Gamma(2\alpha+k)} \,
      P^{(\alpha)}_k(x).
\end{equation}

For Bessel functions, we will denote
\begin{equation}
\label{ec:jad}
    \Jad(x) = 2^\delta \Gamma(\alpha+1)
       \frac{J_{\alpha+\delta+1}(x)}{x^{\alpha+\delta+1}}.
\end{equation}

Also, we will use the Sonine's formula
(see~\cite[\S~12.11, p.~373]{Wat} or~\cite[Lemma~4.13]{StWe})
\begin{equation}
\label{ec:sonine}
    J_{\nu+\mu+1}(z)
    = \frac{z^{\nu+1}}{2^{\nu}\Gamma(\nu+1)}
      \int_0^1 J_{\mu}(zs) s^{\mu+1}(1-s^2)^{\nu} \, ds,
\end{equation}
which is valid for $\Re(\mu)>-1$ and $\Re(\nu)>-1$.

Now, let us see a new expression for the Bochner-Riesz
operator $\BRk$ and its kernel~$\KBRk$. It is more useful for our purposes
than the previous one we showed in~(\ref{ec:KernelBRkintermedia}).

\begin{lem}
\label{lem:Kernel}
Let $\alpha>-1/2$ and $\delta>0$. Then,
\begin{equation*}
    \BRk(f,x)
    = x^{-\alpha} \int_0^{\infty} f(t) \KBRk(x,t) t^{\alpha+1} \, dt
\end{equation*}
with
\begin{equation}
\label{ec:KernelBRk}
    \KBRk(x,t)
    = (xt)^{\alpha} \int_0^{\pi} \Jad(w)
         p^{\alpha}_k(\cos\theta) \sin^{2\alpha}\theta \, d\theta,
\end{equation}
being $w^2 = x^2 + t^2 - 2xt\cos\theta$.
\end{lem}

\begin{proof}
We have already seen that $\BRk(f,x)$ can be written with $\KBRk(x,t)$ as
in~(\ref{ec:KernelBRkintermedia}). Now, let us transform this expression to
become~(\ref{ec:KernelBRk}). For this, we will use
\begin{equation*}
    \int_0^{\pi}\frac{J_{\alpha}(w)}{w^{\alpha}} \,
      p^{\alpha}_k(\cos\theta) \sin^{2\alpha}\theta \, d\theta
    = \frac{J_{\alpha+k}(x)}{x^\alpha} \frac{J_{\alpha+k}(t)}{t^{\alpha}},
\end{equation*}
with $p_k^\alpha$ as in~(\ref{ec:ultraesf}); this formula, that
is valid for $\alpha>-1/2$, can be found in~\cite[\S~11.41, p.~367]{Wat}.
In this way,
\begin{align*}
    K_{\alpha+k}^{\delta}(x,t)
      &= \int_0^1 s(1-s^2)^\delta J_{\alpha+k}(xs)J_{\alpha+k}(ts)\,ds \\
      &= (xt)^\alpha
        \int_0^1 s^{2\alpha+1}(1-s^2)^{\delta}
          \left( \int_0^{\pi} \frac{J_{\alpha}(sw)}{(sw)^{\alpha}}
            \, p^{\alpha}_k(\cos\theta)
            \sin^{2\alpha}\theta \, d\theta \right) \, ds.
\end{align*}
Exchanging the order of integration and applying Sonine's
formula~(\ref{ec:sonine}), we can conclude
\begin{align*}
    \KBRk(x,t)
      &= (xt)^\alpha \int_0^{\pi} w^{-\alpha}
         p^{\alpha}_k(\cos\theta) \sin^{2\alpha}\theta
         \left( \int_0^1 s^{\alpha+1}(1-s^2)^{\delta}
            J_{\alpha}(sw) \,ds \right) \, d\theta \\
     &= (xt)^{\alpha}
        \int_0^{\pi} \Jad(w)
            p^{\alpha}_k(\cos\theta) \sin^{2\alpha}\theta \, d\theta.
\end{align*}
\end{proof}

By using standard estimates for ultraspherical polynomials (see, for
instance~\cite[Th.~7.33.1]{Sze}) we have that, for $\alpha \ge 0$,
\begin{equation}
\label{ec:cotaUltraesf}
    |p^\alpha_k(x)| \le 1, \quad x \in [-1,1].
\end{equation}
Also, the well-known estimates for Bessel functions (of order $\nu>-1$)
\begin{equation*}
    |J_\nu(t)| \le C_\nu t^\nu, \ t \in (0,1),
    \qquad
    |J_\nu(t)| \le C_\nu t^{-1/2}, \ t \in (1,\infty),
\end{equation*}
ensure that
\begin{equation}
\label{ec:cotajad}
    \lpnorm{\Jad(x)}{1} = C_{\alpha,\delta} < \infty
\end{equation}
for $\delta > \delta_0 = (2\alpha+1)/2$.

Now, we have all that we will use for

\begin{proof}[Proof of Theorem~\ref{thm:encima}]
By Lemma~\ref{lem:Kernel}, we can write
\begin{equation}
\label{ec:normaBRk}
    \lpnorm{\BRk f}{p} = \biglpnorm{ x^{-\alpha}
         \int_0^{\infty} f(t) \KBRk(x,t) t^{\alpha+1} \, dt }{p}.
\end{equation}
Moreover, by~(\ref{ec:KernelBRk}) and~(\ref{ec:cotaUltraesf}), we get
\begin{equation*}
    \left| \KBRk(x,t) \right|
      \le C \, (xt)^{\alpha} \int_0^{\pi} |\Jad(w)|
         \sin^{2\alpha}\theta \, d\theta.
\end{equation*}
Then, by using this and~(\ref{ec:convo}), we have 
\begin{multline*}
    \qquad
    \left| x^{-\alpha}
      \int_0^{\infty} f(t) \KBRk(x,t) t^{\alpha+1} \, dt
    \right| \\
    \begin{aligned}
      &\le C \int_0^{\infty} |f(t)|
               \left( \int_0^{\pi} |\Jad(w)|
                 \sin^{2\alpha}\theta \,d\theta \right) t^{2\alpha+1} \, dt
                   \qquad \\
      &= C' \int_0^{\infty} |f(t)|\, T^x(|\Jad|,t)
             t^{2\alpha+1}\, dt\\
      &\le C' \, (|\Jad|*|f|)(x).
    \end{aligned}
\end{multline*}
Finally, by applying~(\ref{ec:normaBRk}), (\ref{ec:acotconvo})
and~(\ref{ec:cotajad}), we conclude
\begin{align*}
    \lpnorm{\BRk f}{p}
    & \le C \biglpnorm{\,|\Jad|*|f|\,}{p} \\
    & \le C \lpnorm{\Jad}{1} \lpnorm{f}{p} \\
    & \le C' \lpnorm{f}{p}.
\end{align*}
\end{proof}


%%%%%%%%%%%%%%%%%%%%%%%%%%%%%%%%%%%%%%%%%%%%%%%%%%%%%%%%%%%%%%%%%%%%%%%%%%%
\section{Uniform boundedness for $\delta$ below the critical index}
\label{sec:debajo}
%%%%%%%%%%%%%%%%%%%%%%%%%%%%%%%%%%%%%%%%%%%%%%%%%%%%%%%%%%%%%%%%%%%%%%%%%%%

The proof of Theorem~\ref{thm:debajo} needs some prior machinery.
In particular, we will use Stein's theorem on interpolation of analytic
families of operators; it can be seen in~\cite{St-Int}. Here, we will show
this result adapted to our spaces~$L^p(x^{2\alpha+1})$.

First, let us consider the notion of analytic family of operators. A family of
operators $\{T_z\}$ depending on a complex parameter $z$ that runs in
$0 \le \Re(z) \le 1$ is called analytic if:

(a) For each $z$, $T_z$ is a linear transformation of simple functions on
$(0,\infty)$ into measurable functions on $(0,\infty)$.

(b) If $\phi$ and $\psi$ are simple functions on $(0,\infty)$, then
\begin{equation*}
    \Phi(z) = \int_0^\infty T_z(\psi,x)\phi(x) \, dx
\end{equation*}
is analytic in $0 < \Re(z) < 1$ and continuous in $0 \le \Re(z) \le 1$.

We say that an analytic family $\{T_z\}$ is of admissible growth if $\Phi(z)$
is of admissible growth; that is, if
\begin{equation*}
    \sup_{|y|\le s} \sup_{0 \le x \le 1} \log |\Phi(x+iy)| \le Ae^{as},
\end{equation*}
where $a<\pi$ and $A$ is a constant. Both $A$ and $a$ may depend on the
functions $\phi$ and~$\psi$.

The interpolation result is therefore:
\begin{thm}
\label{thm:interpolacion}
Let $\{T_z\}$ be an analytic family of linear operators of admissible growth,
defined in the strip $0 \le \Re(z) \le 1$. Suppose that
$1 \le p_0,p_1,q_0,q_1 \le \infty$ and
\begin{equation*}
    \frac{1}{p} = \frac{1-r}{p_0} + \frac{r}{p_1},
    \quad
    \frac{1}{q} = \frac{1-r}{q_0} + \frac{r}{q_1}
\end{equation*}
for $0 \le r \le 1$. We also assume
\begin{equation}
\label{ec:cotasTiy}
    \lpnorm{T_{iy}f}{q_0} \le A_0(y) \lpnorm{f}{p_0}
\end{equation}
and
\begin{equation}
\label{ec:cotasT1+iy}
    \lpnorm{T_{1+iy}f}{q_1} \le A_1(y) \lpnorm{f}{p_1}
\end{equation}
for any simple function $f$. Finally, suppose
\begin{equation*}
%\label{ec:cotasAi}
    \log{|A_i(y)|} \le Ae^{a|y|},\ a<\pi, \quad i = 0,1.
\end{equation*}
Then, we may conclude
\begin{equation}
\label{ec:tesisStein}
    \lpnorm{T_r f}{q} \le A_r \lpnorm{f}{p}
\end{equation}
where
\begin{equation*}
    \log{A_r} \le \int_{\R} \omega(1-r,y) \log{A_0(y)} \, dy
      + \int_{\R} \omega(r,y) \log{A_1(y)} \, dy
\end{equation*}
and
\begin{equation*}
    \omega(1-r,y) = \frac{\tan(\pi r/2)}%
          {2 [\tan^2(\pi r/2)+\tanh^2(\pi y/2)] \cos^2(\pi y/2)}.
\end{equation*}
\end{thm}

\begin{proof}
See~\cite{St-Int} or~\cite[Ch.~V, \S~4]{StWe}.
\end{proof}

\begin{rem}
Note that, if the family $\{T_z\}$ depends on a parameter~$k$, and the
estimates~(\ref{ec:cotasTiy}) and~(\ref{ec:cotasT1+iy}) are independent
of~$k$, we can conclude that the boundedness~(\ref{ec:tesisStein}) will be
uniform in~$k$.
\end{rem}

The proof of Theorem~\ref{thm:debajo} (the uniform boundedness of $\BRk$
in the range $0<\delta\le\delta_0$) will use an analytic family of operators
related with~$\BRk$.

Let us consider $\delta(z) = (1-z)\delta_0 + \varepsilon$ with $\varepsilon>0$,
$0 \le \Re(z) \le 1$ and $\delta_0 = (2\alpha+1)/2$.
We will take the family of operators
\begin{equation}
\label{ec:deltaz}
    \BRkd{z}(f,x)
    = x^{-\alpha} \int_0^{\infty} f(t) \KBRkd{z}(x,t) t^{\alpha+1} \, dt,
    \quad k\in\N,
\end{equation}
where the kernel $\KBRkd{z}$ is as in~(\ref{ec:KernelBRk}) with $\delta$
changed by $\delta(z)$. This definition of $\BRkd{z}$ is valid for 
simple functions in $(0,\infty)$.

Bessel functions of complex order $\lambda=\nu+i\mu$ satisfy
\begin{equation*}
    J_\lambda(t) = \frac{(t/2)\lambda}{\Gamma(1/2)\Gamma(\lambda+1/2)}
      \int_0^1 (1-s^2)^{\lambda-1/2}\cos(st) \, ds, \quad \nu > -1/2
\end{equation*}
(it is just a particular case of~(\ref{ec:sonine})), and the estimates
\begin{gather}
    \label{ec:cota1/2}
    |J_{\nu+i\mu}(t)| \le C_\nu e^{\pi|\mu|} t^{-1/2},
    \quad t \ge 1,\ \nu \ge 0, 
    \\
    \label{ec:cotanu}
    |J_{\nu+i\mu}(t)| \le C_\nu e^{\pi|\mu|} t^{\nu},
    \quad t > 0,\ \nu \ge 0.
\end{gather}
Then, it is not difficult to check that, for any $k$, the family of
operators~(\ref{ec:deltaz}) is analytic and of admissible growth.
The details can be seen in~\cite[\S~3]{Well}.

Now, let us prove the uniform boundedness for the operators $\BRkd{z}$ in
$\Re(z)=0$ and $\Re(z)=1$ for some values of~$p$. 
Lemma~\ref{lem:cotaiy} will establish the boundedness for $\Re(z)=0$,
and Lemma~\ref{lem:cota1+iy} for $\Re(z)=1$.

\begin{lem}
\label{lem:cotaiy}
Let $\alpha\ge 0$ and $1 \le p \le \infty$.
Consider $\BRkd{z}$ the family of operators
given by~(\ref{ec:deltaz}), where $\delta(z) = (1-z)\delta_0+\varepsilon$ with
$\varepsilon>0$, $0 \le \Re(z) \le 1$ and $\delta_0 = (2\alpha+1)/2$.
Then, the inequality
\begin{equation*}
%\label{ec:cotaiy}
    \lpnorm{\BRkd{iy}f}{p} \le A_0(y) \lpnorm{f}{p}, \quad k\in\N,
\end{equation*}
holds, with $A_0(y) = C_{\alpha,p} e^{\pi|\delta_0y|}$.
\end{lem}

\begin{proof}
First, note that Lemma~\ref{lem:Kernel} can be reproduced with the change
$\delta$ by $\delta(iy)$. Then, the proof of this lemma is similar to the one
of Theorem~\ref{thm:encima}, by applying again the convolution structure. But,
this time, we use the estimation
\begin{equation*}
    \lpnorm{\Jad[\delta(iy)](x)}{1}
    \le C_\alpha e^{\pi|\delta_0y|},
\end{equation*}
that follows from the bounds~(\ref{ec:cota1/2}) and~(\ref{ec:cotanu})
for~(\ref{ec:jad}) (with $\delta(iy)$).
\end{proof}

Now, the boundedness for $\BRkd{1+iy}$:

\begin{lem}
\label{lem:cota1+iy}
Let $\alpha\ge 0$ and $1 \le p \le \infty$.
Consider $\BRkd{z}$ the family of operators given by~(\ref{ec:deltaz}), where
$\delta(z) = (1-z)\delta_0+\varepsilon$ with
$\varepsilon>0$, $0 \le \Re(z) \le 1$ and $\delta_0 = (2\alpha+1)/2$. If
\begin{equation}
\label{ec:condiciones}
    \frac{4(\alpha+1)}{2\alpha+3} < p < \frac{4(\alpha+1)}{2\alpha+1},
\end{equation}
then the inequality
\begin{equation*}
%\label{ec:cota1+iy}
    \lpnorm{\BRkd{1+iy}f}{p} \le A_1(y) \lpnorm{f}{p}, \quad k\in\N,
\end{equation*}
holds with
$A_1(y) = C_{\alpha,p} \big(1+\frac{|\delta_0y|}{\varepsilon}\big)$.
\end{lem}

\begin{proof}
Von Lommel's formula
\begin{equation*}
    sJ_{\nu}(xs)J_{\nu}(ts) = \frac{1}{x^2-t^2}
      \frac{d}{ds} \{stJ_{\nu}(xs)J'_{\nu}(ts) - sxJ'_{\nu}(xs)J_{\nu}(ts)\}
\end{equation*}
and an integration by parts in~(\ref{ec:KernelBRkintermedia}) (with $\delta$
changed by $\delta(1+iy)$) give
\begin{equation*}
    \KBRkd{1+iy}(x,t)
    = 2 \delta(1+iy) \int_0^1 s(1-s^2)^{\delta(1+iy)-1} \calK(x,t) \, ds,
\end{equation*}
where
\begin{equation}
\label{ec:calK}
    \calK[\nu](x,t)
    = \frac{stJ_{\nu}(xs)J'_{\nu}(ts)
        - sxJ'_{\nu}(xs)J_{\nu}(ts)}{x^2-t^2}.
\end{equation}
Also, let us denote
\begin{equation}
\label{ec:calT}
    \calT(f,x)
    = x^{-\alpha} \int_0^{\infty} f(t) \calK(x,t) t^{\alpha+1} \, dt.
\end{equation}
In this way, by using Fubini's theorem, we can write
\begin{equation*}
    \BRkd{1+iy}(f,x)
    = 2 \delta(1+iy) \int_0^1 s(1-s^2)^{\delta(1+iy)-1} \calT(f,x) \, ds.
\end{equation*}

Taking into account that $\delta(1+iy) = \varepsilon-iy\delta_0$ and by
applying Minkowsky's integral inequality, we have
\begin{align*}
    \lpnorm{\BRkd{1+iy}f}{p} 
    & \le C (\varepsilon+|\delta_0y|)
        \biglpnorm{\int_0^1 2s(1-s^2)^{\varepsilon-iy\delta_0-1}
          \calT(f,x) \, ds}{p} \\
    & \le C (\varepsilon+|\delta_0y|) \int_0^1 2s(1-s^2)^{\varepsilon-1}
        \biglpnorm{\calT(f,x)}{p} \, ds.
\end{align*}

Now, we claim that
\begin{equation*}
    \lpnorm{\calT f}{p}
    \le C_{\alpha,p} \lpnorm{f}{p}
\end{equation*}
holds for every $s$ for $p$ verifying~(\ref{ec:condiciones}).
We have
$\calT(f(\cdot),x) = \calT[1](f(\frac{(\cdot)}s),sx)$,
and then it is enough to prove the inequality
\begin{equation}
\label{ec:cotacalT1}
    \lpnorm{\calT[1]f}{p} \le C_{\alpha,p} \lpnorm{f}{p}.
\end{equation}
But comparing~(\ref{ec:calK}) and~(\ref{ec:calT}) with~(\ref{ec:makJJ'}), it
is clear that
\begin{equation*}
    \calT[1]f = \mak f.
\end{equation*}
Then, by applying the uniform boundedness of Theorem~\ref{thm:mak}, we
get~(\ref{ec:cotacalT1}).

Finally, using that
\begin{equation*}
    (\varepsilon+|\delta_0y|) \int_0^1 2s(1-s^2)^{\varepsilon-1} \, ds
    \le \Big( 1+\frac{|\delta_0y|}{\varepsilon} \Big),
\end{equation*}
we conclude
\begin{equation*}
    \lpnorm{\BRkd{1+iy}f}{p}
    \le C_{\alpha,p} \Big( 1+\frac{|\delta_0y|}{\varepsilon} \Big)
    \lpnorm{f}{p}.
\end{equation*}
\end{proof}

Now,

\begin{proof}[Proof of Theorem~\ref{thm:debajo}]
We will restrict our attention to $1\le p \le 2$. The other values for $p$
can be obtained by using a duality argument, because
\begin{equation*}
    \int_0^\infty f(x) \BRk(g(t),x) x^{2\alpha+1} \, dx
    = \int_0^\infty \BRk(f(x),t) g(t) t^{2\alpha+1} \, dt.
\end{equation*}

To prove that~(\ref{ec:condprin}) is a necessary condition, let us take
a function $f \in S^+$ such that $\htr(f,x) = 1$ if $x \in (0,1)$. Thus,
by using the definition~(\ref{ec:BRk}) and Sonine's formula~(\ref{ec:sonine}),
we get that
\begin{align*}
    M_{\alpha,0}^{\delta}(f,x)
    &= \htr((1-t^2)^\delta_+ \htr f, x) = \htr((1-t^2)^\delta_+, x) \\
    &= \int_0^1 \frac{J_\alpha(xt)}{(xt)^\alpha} \,
          (1-t^2)^\delta t^{2\alpha+1} \, dt
      = 2^\delta \Gamma(\delta+1)
          \frac{J_{\alpha+\delta+1}(x)}{x^{\alpha+\delta+1}}.
\end{align*}
Now, by applying the well-known asymptotic estimate (see~\cite[\S~7.21
(1), p.~199]{Wat})
\begin{equation*}
    J_\nu(x) \approx {(\tfrac2\pi)}^{1/2} x^{-1/2}
       \cos(x-\tfrac{\nu\pi}2-\tfrac\pi4) + O(x^{-3/2}),
    \quad x\to\infty,
\end{equation*}
it is clear that $M_{\alpha,0}^{\delta}f \notin L^p(x^{2\alpha+1})$
if $p \le \frac{4(\alpha+1)}{2\alpha+2\delta+3}$.

Finally, let us prove that~(\ref{ec:condprin}) is also sufficient.
Let us take the analytic family given by~(\ref{ec:deltaz}), with
$\delta(z) = (1-z)\delta_0+\varepsilon$,
$\delta_0 = (2\alpha+1)/2$ and $\varepsilon>0$,
and take $p_0 = q_0 = 1$, $p_1 = q_1 = 4(\alpha+1)/(2\alpha+3)+\varepsilon$.
Now, let us use Lemmas~\ref{lem:cotaiy} and~\ref{lem:cota1+iy} for
$p_0$ and $p_1$. Then,
\begin{equation*}
   \lpnorm{\BRkd{iy}f}{p_0} \le A_0(y)\lpnorm{f}{p_0},
\end{equation*}
with $A_0(y) = C_{\alpha,p_0} e^{\pi|\delta_0y|}$,
and
\begin{equation*}
    \lpnorm{\BRkd{1+iy}f}{p_1} \le A_1(y) \lpnorm{f}{p_1},
\end{equation*}
with
$A_1(y) = C_{\alpha,p_1} \big(1+\frac{|\delta_0y|}{\varepsilon}\big)$.
In this way, Theorem~\ref{thm:interpolacion} ensures
\begin{equation}
\label{ec:ultima}
    \lpnorm{\BRkd{t}f}p \le C \lpnorm{f}p
\end{equation}
for those values for $p$ which satisfy
\begin{equation*}
    \frac1p = (1-r) + \frac{r}{p_1}, \quad 0 \le r \le 1.
\end{equation*}
By using $\delta = \delta(r) = (1-r)\delta_0 + \varepsilon$, we have
$r = 1 - \frac{\delta-\varepsilon}{\delta_0}$, so the result holds for those
values for $p$ which satisfy
\begin{equation*}
    \frac1p = \Big( \frac{\delta-\varepsilon}{\delta_0} \Big)
      \Big( 1-\frac1{p_1} \Big) + \frac{1}{p_1},
    \quad \varepsilon \le \delta \le \varepsilon+\delta_0.
\end{equation*}
Taking an arbitrarily small $\varepsilon$, we see that, if $0 < \delta <
\delta_0$, the boundedness~(\ref{ec:ultima}) holds for those values of $p$
which satisfy
$\frac1p < \frac{2\alpha+3+2\delta}{4(\alpha+1)}$. This completes the proof.
\end{proof}

%%%%%%%%%%%%%%%%%%%%%%%%%%%%%%%%%%%%%%%%%%%%%%%%%%%%%%%%%%%%%%%%%%%%%%%%%%%
% References
%%%%%%%%%%%%%%%%%%%%%%%%%%%%%%%%%%%%%%%%%%%%%%%%%%%%%%%%%%%%%%%%%%%%%%%%%%%

\begin{thebibliography}{99}

\bibitem{BeRo} J. J. Betancor and L. Rodr\'{\i}guez-Mesa,
$L_1$-convergence and strong summability of Hankel transforms,
\textit{Publ. Math. Debrecen}
\textbf{55} (1999), 437--451.

\bibitem{BeStem} J. J. Betancor and K. Stempak,
Relating multipliers and transplantation for Fourier-Bessel expansions and
Hankel transform,
\textit{T\^ohoku Math. J.}
\textbf{53} (2001), 109--129.

\bibitem{Cho} F. M. Cholewinski,
A Hankel convolution complex inversion theory, 
\textit{Mem. Amer. Math. Soc.}
\textbf{58} (1965).

\bibitem{CGPV-FN} \'O. Ciaurri, J. J. Guadalupe, M. P\'erez and J. L. Varona,
{Mean and almost everywhere convergence of Fourier-Neumann series},
\textit{J. Math. Anal. Appl.}
\textbf{236} (1999), 125--147.

\bibitem{CGPV-DIE} \'O. Ciaurri, J. J. Guadalupe, M. P\'erez and J. L. Varona,
{Solving dual integral equations on Lebesgue spaces},
\textit{Studia Math.}
\textbf{142} (2000), 253--267.

\bibitem{ColCTV} L. Colzani, A. Crespi, G. Travaglini and M. Vignati,
{Equiconvergence theorems for Fourier-Bessel expansions with applications to
the harmonic analysis of radial functions in euclidean and noneuclidean
spaces},
\textit{Trans. Amer. Math. Soc.}
\textbf{338} (1993), 43--55.

\bibitem{Dur1} A. J. Dur\'an,
{Laguerre expansions of tempered distributions and generalized functions},
\textit{J. Math. Anal. Appl.}
\textbf{150} (1990), 166--180.

\bibitem{Dur2} A. J. Dur\'an,
{On Hankel transform},
\textit{Proc. Amer. Math. Soc.}
\textbf{110} (1990), 417--424.

\bibitem{GasTr} G. Gasper and W. Trebels,
{Necessary conditions for Hankel multipliers},
\textit{Indiana Univ. Math. J.}
\textbf{31} (1982), 403--414.

\bibitem{GPV-Con} J. J. Guadalupe, M. P\'erez and J. L. Varona,
{Commutators and analytic dependence of Fourier-Bessel series on $(0,\infty)$},
\textit{Canad. Math. Bull.}
\textbf{42} (2) (1999), 198--208.

\bibitem{Her} C. S. Herz,
{On the mean inversion of Fourier and Hankel transforms},
\textit{Proc. Nat. Acad. Sci. U.S.A.}
\textbf{40} (1954), 996--999.

\bibitem{Hai} D. T. Haimo,
{Integral equations associated with Hankel convolutions},
\textit{Trans. Amer. Math. Soc.}
\textbf{116} (1965), 330--375. 

\bibitem{Hir} I. I. Hirschman, Jr.,
{Variation diminishing Hankel transforms},
\textit{J. Analyse Math.}
\textbf{8} (1960/61), 307--336.

\bibitem{Iga} S. Igari,
{On the multipliers of Hankel transform},
\textit{T\^ohoku Math. J.}
\textbf{24} (1972), 201--206.

\bibitem{Men} J. M. M\'endez,
{On the Bessel transformation of arbitrary order},
\textit{Math. Machr.}
\textbf{136} (1988), 233--239.

\bibitem{RF} J. L. Rubio de Francia,
{Transference principles for radial multipliers},
\textit{Duke Math. J.}
\textbf{58} (1989), 1--19.

\bibitem{St-Int} E. M. Stein,
{Interpolation of linear operators},
\textit{Trans. Amer. Math. Soc.}
\textbf{83} (1956), 482--492.

\bibitem{St-HA} E. M. Stein,
\textit{``Harmonic Analysis: Real-Variable Methods, Orthogonality and
Oscillatory Integrals,''}
Princeton University Press, Princeton (New Jersey), 1993.

\bibitem{StWe} E. M. Stein and G. Weiss,
\textit{``Introduction to Fourier Analysis on Euclidean Spaces,''}
Princeton University Press, Princeton (New Jersey), 1975.
% (es la segunda impresi\'on de la primera edici\'on de 1971)

\bibitem{StemTr} K. Stempak and W. Trebels,
{Hankel multipliers and transplantation operators},
\textit{Studia Math.}
\textbf{126} (1997), 51--66.

\bibitem{Sze} G. Szeg\H{o},
\textit{``Orthogonal Polynomials''} (4th edition),
Amer. Math. Soc. Colloq. Publ.
\textbf{23}, Amer. Math. Soc.,
Providence, RI, 1975.

\bibitem{Var-Evora} J. L. Varona,
{Acotaci\'on en norma de series de Fourier-Neumann},
in ``Actas XV Jornadas Luso-Espanholas de Matem\'atica,'' Vol.~II
(\'Evora, Portugal, 1990), 119--124,
Universidade de \'Evora, \'Evora, 1991.

\bibitem{Var} J. L. Varona,
{Fourier series of functions whose Hankel transform is supported on
$[0,1]$},
\textit{Constr. Approx.}
\textbf{10} (1994), 65--75.

\bibitem{Wat} G. N. Watson,
\textit{``A Treatise on the Theory of Bessel Functions,''}
Cambridge Univ. Press, 1966.

\bibitem{Well} G. V. Welland,
{Mean convergence of Riesz-Bochner means for radial functions},
\textit{Can. J. Math.}
\textbf{27} (1975), 176--185.

\end{thebibliography}

\end{document}
